\documentclass[10pt,a4paper]{amsart}
\usepackage[margin=19mm]{geometry}
\usepackage{amsmath,amssymb,mathtools}
\usepackage[scr=rsfs,cal=euler]{mathalfa}
\usepackage{xfrac}
\usepackage[unicode,hidelinks,bookmarks=false]{hyperref}
\newcommand{\ie}{i.e.\ }
\newcommand{\cf}{cf.\ }
\newcommand{\resp}{respectively\ }
\newcommand{\Subsection}{\subsection}
\providecommand{\nobreakdash}{\nobreak\hbox{-}}
\setlength{\emergencystretch}{2em}
\multlinegap0pt
\usepackage{amssymb}
\usepackage{mathtools}
\usepackage{cancel}
\usepackage[shortlabels]{enumitem}
\usepackage{dsfont}
\newtheorem{lemma}{Lemma}
\newcommand{\C}{\mathbb{C}}
\newcommand{\End}{\operatorname{End}}
\newcommand{\R}{\mathbb{R}}
\newcommand{\pa}{\partial}
\title{The conformal limit for Nakajima quiver varieties}
\author{Sotiria Chatzimarkou, Panagiotis Dimakis}
\date{Source: arXiv:2604.17006v2 (CC BY 4.0)}
\begin{document}
\maketitle

\noindent\emph{Source and licence.} The conformal limit for Nakajima quiver varieties. Version arXiv:2604.17006v2 (CC BY 4.0), distributed by arXiv under the Creative Commons Attribution 4.0 International licence, http://creativecommons.org/licenses/by/4.0/, as recorded in the arXiv licence metadata for this exact version and shown on the abstract page https://arxiv.org/abs/2604.17006v2. Full licence notice: \texttt{licenses/arxiv-2604.17006-CC-BY-4.0.txt}.\par\smallskip

\noindent\textit{Editorial scope.} Counted unit: the article's lemma metric (source0.tex lines 1310--1324). The article's complete original proof of it is delivered with it (source0.tex lines 1394--1411, one \texttt{proof} environment of 18 lines, which the article places away from the statement). No mathematics is written, repaired or re-derived: the statement and the proof are the article's own lines, in the article's own notation, and no retained line is substituted or re-flowed.  Retained uncounted as the source-local context the proof needs: lines 587--589.  Nothing was omitted from any retained span, so the package carries no omission record.  No retained line contains a citation, so the wrapper carries no bibliography and no bibliography entry.  No retained line contains a figure environment, an image-inclusion command or drawing code, so no image is delivered and the package declares no asset dependency.  Editorial formatting only, in addition: an AMS \texttt{amsart} preamble that restates the article's own declarations and macros used by the retained lines, namely the environments \texttt{lemma} on separate counters, together with the standard packages those lines need.  The retained environments print in this document as Lemma 1 in order of appearance, whereas the article prints the same items with the numbering of its own full document.  The complete native source of the article as distributed in the arXiv e-print for this exact version is bundled unchanged as \texttt{sources/198712/source0.tex}.
\begin{equation}\label{CSA}
    \xi\cdot [(\lambda,\bar\pa_E,\nabla_{\lambda})]:= [(\xi\lambda,\bar\pa_E,\xi\nabla_{\lambda})],
\end{equation}

\begin{lemma}\label{metric}
Let $A= (A_h,I_k,A_{\bar h}, J_k) \in S_{p_0}^+$ and write $(B_h,i_k,B_{\bar h},j_k) = p_0 +A_R$. Then there exists $R_A>0$ sufficiently small such that for all $R\le R_A$ there exists a $\xi_R \in i\mathfrak g_{\mathbf{v}}$ which satisfies the real moment map equation 
%
\begin{equation}
\mu_{\R}(e^{\xi_R}\cdot(p_0+A_R))_k = \zeta_{\R}^k
\end{equation}
%
and 
%
\begin{equation}\label{metricansatz}
(\xi_R)_k \in \End(V_k)_0(R^2) + \bigoplus\limits_{|m|= 1}\End(V_k)(R^3)+ ... + \bigoplus\limits_{|m|= M} \End(V_k)_m(R^{M+2})
\end{equation}
%
with $M$ being the maximal positive integer for which $\End(V_k)_M \neq \emptyset$.
\end{lemma}

\begin{proof}
We already know this when the norm of $A$ is small by Lemma \ref{metric} above. Therefore assume that we know that $[p_0+A]\in W_a^0(p_0)$. Now by Lemma \ref{CSA} and the definition of $W_a^0(p_0)$ we know that this implies that $(B_h^0+A_h,i_k^0+I_k,R(B_{\bar h}^0+A_{\bar h}), R(j_k^0+J_k))$ is also in the complex gauge orbit of a point in $W_a^0(p_0)$. This implies that there exists a unique $g^1_A$ such that 
%
\begin{equation*}
\mu_{\R}(g^1_A\cdot(B_h^0+A_h,i_k^0+I_k,R(B_{\bar h}^0+A_{\bar h}), R(j_k^0+J_k))) = \zeta_{\R}.
\end{equation*}
%
But then for $g_A = g^1_Ag^{-1}_R$ it holds that 
%
\begin{equation*}
\begin{split}
\zeta_{\R} &= \mu_{\R}(g^1_A\cdot(B_h^0+A_h,i_k^0+I_k,R(B_{\bar h}^0+A_{\bar h}), R(j_k^0+J_k))) \\
&= \mu_{\R}(g_A\cdot(p_0+A_R)).
\end{split}
\end{equation*}
%
Therefore, once the theorem hods for small $A\in S^+_{p_0}$ we can use the $\C^{\star}$-action to scale it to any $A$. 
\end{proof}

\end{document}
